% =====================================================================
%   HGND + GNN sensitivity paper — DRAFT SKELETON
%   Target: Phys. Rev. C (or equivalent Q1: NIM A, EPJ C, JHEP)
%   Style : revtex4-2, twocolumn, superscript numeric refs (aps style)
%
%   Structure and voice inspired by:
%     Morozov et al., "The Highly-Granular Time-of-Flight Neutron
%     Detector for the BM@N experiment", arXiv:2412.00455v1 (Nov 2024)
%   Latest ML/physics results synthesised from the presentation:
%     V. Bocharnikov, "Graph Neural Network-based neutron reconstruction
%     in the HGND at the BM@N experiment", Analysis and Detector Meeting
%     of the BM@N, 26.02.2026 (docs/VB_BMN_Dombay.pdf).
%
%   Fill-in policy
%   --------------
%   Every unresolved section is marked
%       %% TODO(<tag>): <what is missing>  [<owner>]
%   Every quantitative claim that needs verification is
%       \NUM{value}{units}{tag}    ← replace with the measured number
%   Every citation that needs a bib entry is
%       \CITENEEDED{short-tag}
%   Search this file for the string TODO / NUM / CITENEEDED before
%   submitting.  Nothing should remain marked.
%
%   Sensitivity study is the *headline* result.  Until the full-scale
%   SLURM sweep produces the final N_true(E_k) curves and closure
%   diagnostic, the section carries placeholders + explicit blockers.
% =====================================================================

\documentclass[aps,prc,twocolumn,superscriptaddress,showpacs,%
               nofootinbib,floatfix,longbibliography]{revtex4-2}

\usepackage[utf8]{inputenc}
\usepackage[T1]{fontenc}
\usepackage{amsmath,amssymb,amsfonts}
\usepackage{graphicx}
\usepackage{xcolor}
\usepackage{siunitx}
\usepackage{hyperref}
\usepackage{booktabs}
\usepackage{physics}
\hypersetup{colorlinks=true,linkcolor=blue,citecolor=blue,urlcolor=blue}

% ---- Editorial helpers (remove before submission) --------------------
\newcommand{\TODO}[1]{{\color{red}\textbf{[TODO: #1]}}}
\newcommand{\NOTE}[1]{{\color{orange}\textit{[note: #1]}}}
\newcommand{\NUM}[3]{{\color{purple}\underline{#1}~#2}\,\NOTE{#3}}
\newcommand{\CITENEEDED}[1]{{\color{red}\textbf{[cite: #1]}}}

% ---- Physics shorthand ----------------------------------------------
\newcommand{\Ekin}{E_{\text{kin}}}
\newcommand{\Nreco}{N_{\text{reco}}}
\newcommand{\Ntrue}{N_{\text{true}}}
\newcommand{\Nmc}{N_{\text{MC}}}
\newcommand{\eff}{\varepsilon}
\newcommand{\Usym}{U_{\text{sym}}}
\newcommand{\NperEv}[1]{#1/N_{\text{ev}}}
\newcommand{\HGND}{HGND}
\newcommand{\BMN}{BM@N}
\newcommand{\SMASH}{\textsc{smash}}
\newcommand{\DCM}{\textsc{dcm-qgsm-smm}}
\newcommand{\GEANT}{\textsc{Geant4}}
\newcommand{\PyG}{PyTorch~Geometric}

% =====================================================================
\begin{document}

\title{Sensitivity of the Highly-Granular Time-of-Flight Neutron
       Detector at \BMN{} to the nuclear symmetry potential\\
       via graph-neural-network based neutron reconstruction}

% -- Author list: mirror the arXiv paper 2412.00455v1 ordering; ML
%    contributors bumped forward for this ML/physics-driven paper.
\author{V.~Bocharnikov}
\email{vbocharnikov@hse.ru}
\affiliation{HSE University, Moscow, Russia}
\author{S.~Morozov}
\affiliation{Institute for Nuclear Research of the Russian Academy of
             Sciences, 60-letiya Oktyabrya prospekt 7a, 117312 Moscow,
             Russia}
\author{D.~Finogeev}
\author{M.~Golubeva}
\author{F.~Guber}
\author{A.~Izvestnyy}
\author{N.~Karpushkin}
\author{A.~Makhnev}
\author{P.~Parfenov}
\author{D.~Serebryakov}
\author{A.~Shabanov}
\author{V.~Volkov}
\author{A.~Zubankov}
\affiliation{Institute for Nuclear Research of the Russian Academy of
             Sciences, 60-letiya Oktyabrya prospekt 7a, 117312 Moscow,
             Russia}
%% TODO(authors): confirm final author list + order + affiliations
%%                with the BM@N HGND working group before submission.

\date{\today}

\begin{abstract}
%% TODO(abstract): rewrite once the sensitivity result is final.
%% Draft:
We investigate the sensitivity of the Highly-Granular Time-of-Flight
Neutron Detector (\HGND{}) of the fixed-target \BMN{} experiment at
JINR to the nuclear symmetry potential $\Usym$ using a graph
neural-network (GNN) based neutron reconstruction algorithm.  The
detector concept, geometry, and time-of-flight response are those of
Morozov \emph{et~al.}~\cite{Morozov2024HGND}.
Neutron identification, energy regression, and multiplicity
reconstruction are performed cluster-wise with a heterogeneous GNN
trained on \DCM{} + \GEANT{} Bi+Bi and \SMASH{} Xe+CsI simulations.
The reconstructed per-event neutron kinetic-energy spectra
$\Nreco(\Ekin)$ are efficiency-corrected against a Monte-Carlo (MC)
truth benchmark to yield the solved neutron yield
$\Ntrue(\Ekin) = \Nreco(\Ekin)/\eff(\Ekin)$.
We compare $\Ntrue(\Ekin)$ across three \SMASH{} datasets with symmetry
potentials $\Usym = \{0, 18, 90\}$~MeV
(\texttt{zeroSpot}/\texttt{defaultSpot}/\texttt{bigSpot}) and quantify
the pipeline's sensitivity to $\Usym$ as a function of $\Ekin$, decision
threshold, and reconstruction model.
Key numerical results:
reconstruction efficiency
$\langle\eff\rangle \approx \NUM{80}{\%}{smoke\_placeholder}$ at
purity $\approx \NUM{87}{\%}{smoke\_placeholder}$;
energy resolution below 10\% for $0.7 \lesssim \Ekin \lesssim 5$~GeV;
cross-dataset yield ratios \TODO{quote from final SLURM sweep with
propagated uncertainties}.
\end{abstract}

\maketitle

% =====================================================================
\section{Introduction}
\label{sec:intro}
% ---------------------------------------------------------------------
% Following 2412.00455v1 §1: EoS motivation → symmetry-energy gap →
% BM@N/HGND capability → gap this paper fills (ML + cross-dataset
% sensitivity). ~1.5 pages in single-column proof.
% ---------------------------------------------------------------------

The equation of state (EoS) of dense nuclear matter, defined as the
binding energy per nucleon as a function of pressure, density, energy
and temperature, separates into a symmetric-matter component and a
symmetry-energy component~\cite{Burgio2020NuclearMatter}.  The
symmetric-matter term is well constrained by heavy-ion measurements of
meson yields and anisotropic flow at intermediate
energies~\cite{E895_ellipticFlow, E895_sideFlow, E895_diffElliptic,
Sorensen2024EoS}, although residual discrepancies between experiments
translate into a sizeable systematic on the EoS constraints and
motivate new data at higher baryon density~\cite{Sorensen2024EoS,
LeFevre2015}.  The symmetry-energy term encodes the isospin
asymmetry between neutron- and proton-rich matter and is a direct
input to neutron-star observables (radii, tidal deformabilities,
merger dynamics)~\cite{Burgio2020NuclearMatter}.  At high baryon
density the only laboratory constraint comes from the FOPI/LAND
neutron-to-proton anisotropic-flow ratio at
400~MeV/nucleon~\cite{FOPI_LAND1, FOPI_LAND2}; transport-model
studies indicate that neutron yields and $n/p$ flow ratios carry
significant sensitivity to the symmetry potential $\Usym$ in the
$\sim 1$–$4$~AGeV window that has not been experimentally
accessed~\cite{LongWei2022, Sorensen2024EoS}.

The \BMN{} (Baryonic Matter at Nuclotron) fixed-target experiment at
JINR provides ion beams from protons to gold with kinetic energies up
to $6$~AGeV for light ions and $4.5$~AGeV for heavy
ions~\cite{BMN_EPJ, Afanasiev2024BMN}, reaching baryon densities of
3–4~$\rho_0$ in the collision fireball~\cite{Cleymans, Friman}, and
delivers a comprehensive tracking, PID and centrality
setup~\cite{Afanasiev2024BMN, Senger2022BMN}.  First Xe+CsI data were
recorded at $3.8$~AGeV in 2023.  The Highly-Granular Time-of-Flight
Neutron Detector (\HGND{}) has been designed and constructed to
measure neutron yields and azimuthal flow at these energies for
heavier colliding systems~\cite{Morozov2024HGND, Guber2023HGND};
the detector concept, readout, and simulation performance were
reported in Ref.~\cite{Morozov2024HGND}, together with two prototype
neutron reconstruction strategies --- a rule-based cluster method and
a first two-network GNN baseline.  A parallel study established the
intrinsic scintillator time resolution at
$\sim 130$–$150$~ps~\cite{Guber2023SiPM}, and a 100~ps FPGA TDC
readout was demonstrated~\cite{Finogeev2024TDC}.  The estimated
primary neutron flux in the \HGND{} acceptance is
$\sim 500$~neutrons per second in the
$0.3 \lesssim \Ekin \lesssim 4$~GeV range~\cite{Morozov2024HGND}.

The present paper closes two gaps in that pipeline.
First, we replace the two-network baseline (event-level classification
+ energy regression) with a single heterogeneous GNN that operates
directly on the physically motivated hit \emph{clusters}, jointly
producing a per-cluster neutron score and an energy prediction.  This
brings the reconstruction primitives one step closer to the ideal
neutron-per-cluster observable and simplifies downstream analysis.
Second, and more importantly, we quantify the \emph{physics}
sensitivity of the \HGND{} + GNN pipeline to the underlying symmetry
potential by comparing reconstructed neutron kinetic-energy spectra
across three \SMASH{} Xe+CsI datasets differing only in $\Usym$
($0$, $18$, $90$~MeV).  This turns the detector performance study of
Ref.~\cite{Morozov2024HGND} into a physics-observable study, and
delivers the first end-to-end sensitivity benchmark for the \HGND{}
before first data-taking.

The paper is organised as follows.
Section~\ref{sec:detector} briefly recaps the \HGND{} concept for the
reader who does not have Ref.~\cite{Morozov2024HGND} to hand, focussing
on the pieces relevant to reconstruction.
Section~\ref{sec:simulation} describes the two simulation samples
used: \DCM{} + \GEANT{} for detector-response validation and \SMASH{}
Xe+CsI with three $\Usym$ values for the sensitivity study.
Section~\ref{sec:gnn} details the graph construction, network
architecture, and training procedure.
Section~\ref{sec:performance} reports classification, energy, and
multiplicity performance on the \DCM{} sample.
Section~\ref{sec:sensitivity} presents the sensitivity study: MC-truth
spectra, per-event reconstructed yields, the efficiency-corrected
observable $\Ntrue(\Ekin)$, and cross-dataset ratios.
Section~\ref{sec:systematics} discusses systematic uncertainties.
Section~\ref{sec:summary} summarises and outlines the path to first
data.

% =====================================================================
\section{The HGND at BM@N}
\label{sec:detector}
% ---------------------------------------------------------------------
% Compact recap (~0.75 pages) — most of the design detail is deferred
% to 2412.00455v1.  Include only what is needed to read the plots.
% ---------------------------------------------------------------------

The \HGND{} is a sampling calorimeter combining copper absorbers with
high-granularity plastic scintillator readout.  Two geometrical
options were studied in Ref.~\cite{Morozov2024HGND}: a single 16-layer
module at 5~m from the target (17$^\circ$ to the beam axis) and a
split ``2-arms'' option in which two 8-layer modules are placed at
7~m from the target and offset out of the horizontal plane
(10$^\circ$ to the beam axis).  The 2-arms option was adopted to
compensate for the reduction of acceptance at the larger distance and
to keep the rapidity coverage centred around mid-rapidity
($y \sim 1.06$, extending to $y \sim 2.4$) while avoiding
spectator neutrons at forward angles below
$4.7^\circ$~\cite{Morozov2024HGND}.  In this work we use the 2-arms
geometry throughout.

Each arm consists of eight active layers separated by $3$~cm copper
absorbers.  An active layer hosts an $11\times11$ grid of
$4\times4\times2.5$~cm$^3$ polystyrene scintillator cells (with 1.5\%
paraterphenyl + 0.01\% POPOP, produced at
JINR)~\cite{Morozov2024HGND}, packaged in a 3D-printed PETG
light-tight frame capped by 1.5~mm aluminium sheets.  Two downstream
PCBs host $55 + 66$ SiPMs, preamplifiers and LVDS comparators; an
upstream PCB with 121 LEDs enables full-circuit calibration.  A
scintillator veto layer preceding the first active layer rejects
upstream charged particles.  With $121$ cells $\times$ $8$ layers per
arm the pipeline sees $2 \times 968$ instrumented cells.  The
detector-mounted FPGA readout with Kintex~7 TDCs delivers a per-cell
time resolution of $\sim 150$~ps~\cite{Guber2023SiPM,
Finogeev2024TDC}, giving a kinematic reach of
$0.3 \lesssim \Ekin \lesssim 4$~GeV for the time-of-flight (ToF)
neutron energy estimator
\begin{equation}
  E_{\text{ToF}}
    = m_n\!\left(\!\frac{1}{\sqrt{1-\beta^2}} - 1 \right),
  \qquad \beta = \frac{d}{c\,t_{\text{hit}}},
  \label{eq:etof}
\end{equation}
where $d$ and $t_{\text{hit}}$ are the cell-to-vertex distance and
the (noise-broadened) hit time respectively.  The fiducial hit
selection follows Ref.~\cite{Morozov2024HGND}: a per-cell energy
deposition threshold $E_{\text{dep}} > 3$~MeV, corresponding to
approximately $0.5$~MIP, and a hit-time cut $t_{\text{hit}} < 35$~ns
tuned to accept neutrons with $\Ekin > 300$~MeV while suppressing
low-energy background neutrons by a factor of $\sim 6$ at the cost
of $\sim 8\%$ of the primary signal.

\begin{figure}[!t]
  \centering
  \includegraphics[width=\linewidth]{figs/detector_layout_placeholder}
  \caption{
    (Placeholder) Schematic view of the \HGND{} 2-arms geometry at
    \BMN{}.  Reproduced/adapted from Ref.~\cite{Morozov2024HGND}
    Fig.~2 pending permission.
    %% TODO(fig:detector): decide re-use vs redraw for consistent style.
  }
  \label{fig:detector}
\end{figure}

% =====================================================================
\section{Simulation samples}
\label{sec:simulation}
% ---------------------------------------------------------------------
% Two samples: DCM (detector-response validation) + SMASH (sensitivity).
% Report event counts, cuts, and any known differences vs. the arXiv
% paper's DCM sample.
% ---------------------------------------------------------------------

Two independent simulation samples are used in this work.  Detector
response is generated with \GEANT{}~\cite{Geant4} inside the
\texttt{bmnroot}~\cite{bmnroot} simulation and analysis framework;
Gaussian time smearing with $\sigma_t = 150$~ps is applied to each
\HGND{} hit, matching the calibration reported in
Ref.~\cite{Morozov2024HGND}.  The cell-level fiducial selection of
Sec.~\ref{sec:detector} is applied at analysis time.

\paragraph*{DCM sample --- detector-response validation.}
For classification, energy-regression and multiplicity benchmarking we
use $\NUM{300000}{}{Dombay slide 4}$ heavy-ion events generated with
the DCM-QGSM-SMM Monte Carlo cascade
generator~\cite{Baznat2020DCMSMM}, replicating the settings of
Ref.~\cite{Morozov2024HGND} so that the results reported here are
directly comparable to the earlier cluster-method and two-network GNN
baselines.  This sample supports the per-cluster efficiency, purity,
energy-linearity, and multiplicity plots of
Sec.~\ref{sec:performance}.

\paragraph*{SMASH sensitivity sample.}
For the cross-dataset sensitivity study we use three \SMASH{}
\cite{SMASH}
samples of Xe+Cs at $2.87$~AGeV, each containing
$\NUM{500000}{}{Dombay slide 22}$ events, differing only in the
value of the nuclear symmetry potential:
\begin{center}
  \begin{tabular}{@{}lll@{}}
    \toprule
    Dataset label     & $\Usym$ (MeV)  & Preprocessing status \\
    \midrule
    \texttt{zeroSpot}    & $0$   & full CSV dump on cluster            \\
    \texttt{defaultSpot} & $18$  & 200k preprocessed graphs (schema v2)\\
    \texttt{bigSpot}     & $90$  & full CSV dump on cluster            \\
    \bottomrule
  \end{tabular}
\end{center}
Only \texttt{defaultSpot} is fully preprocessed at the time of
writing; the other two are reduced to smoke-scale subsets
(\num{4000} graphs each) for laptop validation.  The full-scale
preprocessing runs on 128~GB nodes and takes
\NUM{4}{\text{h}}{estimate} per dataset (see
Appendix~\ref{sec:reproducibility}).

\paragraph*{Common preprocessing.}
For each event we build a heterogeneous graph as follows.  Cells with
$E_{\text{dep}} > 3$~MeV and $t_{\text{hit}} < 35$~ns become
\texttt{hits} nodes; edges are the intersection of a spatial
radius-graph ($r_{\text{loc}} = 3.6$ cells) and a temporal
radius-graph ($\Delta t = 1.5$~ns).  Connected components of that
graph form \texttt{clusters} nodes; a cluster is labelled ``signal''
if it contains at least one prompt-neutron hit passing the upstream
surface.  Per-cluster energy targets are the fastest prompt neutron's
kinetic energy.  Auxiliary \texttt{events} and \texttt{sides}
node types encode the top/bottom arm assignment for hetero-GNN
compatibility.

% =====================================================================
\section{GNN reconstruction}
\label{sec:gnn}
% ---------------------------------------------------------------------
% Message-passing model + loss + training curriculum.  Diagram helps.
% ---------------------------------------------------------------------

\subsection{Baseline: rule-based cluster reconstruction}
\label{sec:gnn:baseline}

We benchmark against the rule-based cluster reconstruction of
Ref.~\cite{Morozov2024HGND} in a version reimplemented for this
analysis to run on the same preprocessed graphs.  Neighbouring fired
cells with close time stamps are grouped into clusters; the ``head''
of each cluster is the cell with the largest apparent speed
$v = d/t_{\text{hit}}$, and sub-clusters with similar speeds are
merged into larger clusters.  Neutron candidacy is enforced by
rejecting clusters that contain veto-layer or first-active-layer cells
(suppressing charged particles and $\gamma$) and by requiring
$v < c$.  The kinetic energy is obtained from Eq.~\eqref{eq:etof}
applied to the cluster head, and the expected number of fired cells,
sub-clusters and total deposited energy provide a cross-check on the
neutron hypothesis.  This baseline is well understood and has a
predictable systematic behaviour --- notably a positive bias in the
energy determination at $\Ekin > 2$~GeV that we track together with
the GNN prediction throughout the paper.

\subsection{Architecture}

The reconstruction network is a per-cluster classifier + energy
regressor built on \PyG{}~\cite{FeyLenssen2019PyG,
KipfWelling2016GCN}.  The hit branch is
a stack of \texttt{EdgeConv}~$\to$ \texttt{EdgeConv}~$\to$ $L$
\texttt{SAGEConv} $\to$ \texttt{GraphConv}~$\to$ \texttt{GraphConv}
layers with hidden dimension $H$ (default $H=512$, $L=4$).  Edge
scores from the hit branch are learned via a two-layer MLP and used to
weight message passing.  Hits are pooled into cluster representations
by an \texttt{avg\_pool\_x} operator; the cluster branch is a
\texttt{DynamicEdgeConv} ($k=100$) followed by two MLP heads
producing (i) a sigmoid neutron score $s \in [0,1]$ and (ii) a scalar
energy prediction $E_{\text{pred}}$.  The link-classification head
uses per-hit and per-cluster features to score every event edge.
Total parameter count $\NUM{6.4e6}{}{measure with build\_default\_net}$.

We additionally evaluate three alternative model families:
\texttt{hetero\_sage}, \texttt{hetero\_gat}, and \texttt{hetero\_hgt}
(vkr26 heterogeneous graph attention baselines), and the
spectral-gated \texttt{spectral\_dynedge} of the thesis baseline.
Model-family agreement is one of the systematic checks in
Sec.~\ref{sec:systematics}.

\subsection{Loss}

The composite loss combines four BCE terms and one MSE term:
\begin{equation}
  \mathcal{L}
  = \lambda_h \mathcal{L}^{\text{BCE}}_{\text{hit}}
  + \lambda_l \mathcal{L}^{\text{BCE}}_{\text{link}}
  + \lambda_c \mathcal{L}^{\text{BCE}}_{\text{cl}}
  + \lambda_E \mathcal{L}^{\text{MSE}}_{E,\text{cl}}
  + \lambda_{cl} \mathcal{L}^{\text{BCE}}_{\text{cl-link}}.
  \label{eq:loss}
\end{equation}
Default weights are all unity; ablations vary $\lambda_E$ over
$\{0.1, 0.5, 1, 2\}$.

\subsection{Training procedure}

Training is performed on a single NVIDIA A100 with the Adam optimizer
(learning rate $10^{-3}$, weight decay $10^{-5}$),
\texttt{MultiStepLR} milestones at epochs $\{10, 15, 25\}$
($\gamma=0.1$), batch size 512, and a 50/50 train/test split (seed
42).  Full-scale training runs for 20 epochs and completes in
\NUM{4}{\text{h}}{measure on cluster} wall time.  Model checkpoints,
training curves, and per-run manifests are archived under
\texttt{checkpoints/} of the analysis repository.

\begin{figure}[!t]
  \centering
  \includegraphics[width=\linewidth]{figs/gnn_architecture_placeholder}
  \caption{
    (Placeholder) Schematic of the heterogeneous GNN reconstruction
    network: hit-level message passing, cluster pooling, and two
    per-cluster heads for neutron classification and energy regression.
    %% TODO(fig:arch): produce a clean tikz or PPT diagram.
  }
  \label{fig:arch}
\end{figure}

% =====================================================================
\section{Reconstruction performance}
\label{sec:performance}
% ---------------------------------------------------------------------
% Reproduce and extend the Dombay results on the DCM sample.  These
% panels are copied by convention from NeutronRecoGNN/results_smash.
% ---------------------------------------------------------------------

We benchmark the GNN reconstruction on the \DCM{} Xe+CsI sample
following Ref.~\cite{Morozov2024HGND}, adding cluster-level metrics
that were not available in the earlier baseline.

\subsection{Classification}

Figure~\ref{fig:roc} shows the receiver operating characteristic (ROC)
curves for the hit-level classifier and the cluster-level classifier
(combined top+bottom).  The area under the cluster ROC is
$\text{AUC}_{\text{cl}} \approx \NUM{0.97}{}{Dombay slide 12}$,
consistent with the value reported at the \BMN{}
analysis-and-detector meeting~\cite{Bocharnikov2026Dombay}.  A per-
event operating point at cluster score threshold $t = 0.5$ yields
integrated efficiency $\approx \NUM{0.80}{}{Dombay}$ and purity
$\approx \NUM{0.87}{}{Dombay}$; the full purity/efficiency curve is
overlaid in the right panel.
%% TODO(fig:roc): produce final figure from
%%   notebooks/results/sensitivity_full/performance_defaultSpot.png

\begin{figure}[!t]
  \centering
  \includegraphics[width=\linewidth]{figs/performance_roc_placeholder}
  \caption{Left: hit- and cluster-level ROC curves on the DCM sample
    (blinded to training events).  Right: purity and efficiency
    versus decision threshold for the best-cluster-per-event strategy.
    %% TODO(fig:roc): swap in the final PDF.
  }
  \label{fig:roc}
\end{figure}

\subsection{Energy regression}

Figure~\ref{fig:ereco} shows $E_{\text{reco}}$ vs.\ $E_{\text{true}}$
on a logarithmic colour scale.  Linearity is within 10\% for
$0.7 \lesssim \Ekin \lesssim 5$~GeV, and the model compensates for the
ToF-based overestimate visible at $\Ekin > 2$~GeV in the pure ToF
estimator of Eq.~\eqref{eq:etof} (compare
Ref.~\cite{Morozov2024HGND}, Fig.~13).  The relative energy resolution
is $\sigma_E/E < 10\%$ over the same interval.
%% TODO(sec:performance): quote final resolution vs Ekin curve.

\begin{figure}[!t]
  \centering
  \includegraphics[width=\linewidth]{figs/ereco_vs_etrue_placeholder}
  \caption{Cluster-level $E_{\text{reco}}$ vs. $E_{\text{true}}$ 2D
    histogram (log colour).  Dashed line is $y=x$.  Compare to
    Ref.~\cite{Morozov2024HGND} Fig.~13 (right).
    %% TODO(fig:ereco): swap in final PDF.
  }
  \label{fig:ereco}
\end{figure}

\subsection{Neutron multiplicity}

The multiplicity confusion matrix
$N_{\text{reco}}^{\text{cl}} \times N_{\text{true}}$ at threshold
$t = 0.5$ is shown in Fig.~\ref{fig:mult}.  The pipeline resolves up
to $\NUM{3}{}{Dombay slide 15}$ ``good'' neutron clusters per event;
horizontal (vertical) normalisation quantifies efficiency (purity)
as a function of multiplicity.  The event-level neutron multiplicity
distribution reproduces the shape reported in
Ref.~\cite{Morozov2024HGND} Fig.~9 within statistical uncertainties.
For context, the DCM-QGSM-SMM simulation predicts that single-neutron
events account for $\sim 17.3\%$ of all Bi+Bi @ $3.0$~AGeV events
reaching the \HGND{} surface ($\sim 76\%$ of neutron-containing
events), and $\sim 12.7\%$ for the Xe+CsI @ $3.8$~AGeV
sample~\cite{Morozov2024HGND}.  Folded with the nominal \BMN{}
operation cycle ($10^{6}$~ions/spill, 50\% duty factor, 70\%
Nuclotron efficiency, 2\% target interaction length) and the mean
detection efficiency reported below, Ref.~\cite{Morozov2024HGND}
projects $\sim 1.2$–$1.5\times 10^{9}$ reconstructed single-neutron
events per month of Bi+Bi running and $\sim 0.9\times 10^{9}$ per
month of Xe+CsI running.  These statistics comfortably cover the
multi-dataset sensitivity study of Sec.~\ref{sec:sensitivity}.

\begin{figure}[!t]
  \centering
  \includegraphics[width=\linewidth]{figs/multiplicity_placeholder}
  \caption{Neutron multiplicity: (left) event-level $N_n$ per event,
    (right) confusion matrix $N_{\text{reco}} \times N_{\text{true}}$
    at threshold $t = 0.5$.
    %% TODO(fig:mult): reproduce from Dombay slides 15--16 with the
    %% final SLURM checkpoint.
  }
  \label{fig:mult}
\end{figure}

% =====================================================================
\section{Sensitivity to the nuclear symmetry potential}
\label{sec:sensitivity}
% ---------------------------------------------------------------------
% *** Headline result section ***
% Everything downstream of the SLURM full-scale sweep lands here.
% Current smoke-scale numbers are placeholders; the section outline
% is fixed so full-scale drop-in is mechanical.
% ---------------------------------------------------------------------

The physics observable of interest is the reconstructed neutron
kinetic-energy yield per event as a function of the assumed symmetry
potential.  We compare the three \SMASH{} datasets on equal footing:
identical GNN weights (trained on \texttt{defaultSpot}), identical
number of evaluation events per dataset, identical binning, and
identical decision threshold.

\subsection{MC-truth benchmark}

Before evaluating the reconstruction we plot the per-event MC-truth
neutron spectrum for each dataset, defined as one entry per unique
$(\text{event}, \text{track})$ signal neutron
($\texttt{n0\_label}=1$).  This is the physical spectrum the
pipeline aims to recover in the limit of perfect reconstruction and
zero background, and it also sets the natural upper bound on any
efficiency-corrected reconstructed yield.  Figure~\ref{fig:mctruth}
shows the three spectra.  Any downstream cross-dataset comparison must
either normalise these curves to unit area or explicitly account for
their differences.

\begin{figure}[!t]
  \centering
  \includegraphics[width=\linewidth]{figs/mctruth_spectra_placeholder}
  \caption{Per-event MC-truth neutron kinetic-energy spectra for the
    three \SMASH{} samples (\texttt{zeroSpot}: $\Usym=0$~MeV,
    \texttt{defaultSpot}: $\Usym=18$~MeV, \texttt{bigSpot}:
    $\Usym=90$~MeV).  Error bars are $\sqrt{N}$ Poisson counts.
    %% TODO(fig:mctruth): swap in
    %%   notebooks/results/sensitivity_full/mctruth_spectra.png
  }
  \label{fig:mctruth}
\end{figure}

\subsection{Per-dataset efficiency}

The per-dataset per-bin reconstruction efficiency
\begin{equation}
  \eff(\Ekin, t;\, d)
  = P\!\left( s > t \,\middle|\, y=1,\; E_{\text{true}}\in\text{bin}(\Ekin),\; d\right)
  \label{eq:eps}
\end{equation}
is estimated from the same signal-cluster population that provides the
MC-truth benchmark; error bars are the Wilson binomial score interval
\cite{Wilson1927} so that low-statistics bins do not collapse to
$\{0,1\}$ with vanishing uncertainty.  Figure~\ref{fig:eff}
reports $\eff(\Ekin)$ at the working point $t = 0.5$; the three
datasets differ only through the underlying $\Usym$-dependent MC
spectrum, and any residual $\eff$ mismatch is a systematic that must
be tracked in the closure test below.

\begin{figure}[!t]
  \centering
  \includegraphics[width=\linewidth]{figs/efficiency_vs_ekin_placeholder}
  \caption{Per-dataset reconstruction efficiency $\eff_n(\Ekin)$ at
    $t=0.5$ with Wilson binomial intervals.
    %% TODO(fig:eff): swap in efficiency_vs_ekin.png from full run.
  }
  \label{fig:eff}
\end{figure}

\subsection{Solved neutron yield $\Ntrue(\Ekin)$}

The efficiency-corrected observable is
\begin{equation}
  \Ntrue(\Ekin;\, d)
  = \Nreco(\Ekin;\, d) \,/\, \eff(\Ekin;\, d),
  \label{eq:solved}
\end{equation}
with the uncertainty on $\Ntrue$ propagated in quadrature from the
Poisson error on $\Nreco$ and the Wilson interval on $\eff$.  We
report both raw and efficiency-corrected yields per event
($\NperEv{\Nreco}$ and $\NperEv{\Ntrue}$) and overlay the per-event
MC-truth benchmark as a reference.

\begin{figure}[!t]
  \centering
  \includegraphics[width=\linewidth]{figs/sensitivity_yield_placeholder}
  \caption{Per-event reconstructed yield $\NperEv{\Nreco}$ (left) and
    efficiency-corrected yield $\NperEv{\Ntrue}$ (right) versus
    $\Ekin$ for the three \SMASH{} samples.  Dashed lines are the
    per-event MC-truth benchmark.
    %% TODO(fig:yield): swap in sensitivity_yield_analysis.png.
  }
  \label{fig:yield}
\end{figure}

\subsection{Cross-dataset ratios and closure}

The sensitivity signal is the ratio
$\NperEv{\Ntrue}(d)/\NperEv{\Ntrue}(\text{defaultSpot})$ for
$d \in \{\text{zeroSpot}, \text{bigSpot}\}$.  A pipeline that faithfully
recovers the underlying physics returns ratios that (i) differ from
unity by more than the propagated uncertainties in physically
meaningful $\Ekin$ regions, and (ii) track the corresponding MC-truth
ratios (dashed reference in Fig.~\ref{fig:ratio}).  Any residual
discrepancy between reco and MC-truth ratios is the systematic ceiling
of the current analysis.

We additionally quote a \emph{closure} metric,
$\Ntrue/\Nmc$, integrated over all $\Ekin$ bins.  Values within
$1\pm\sigma_{\text{stat}}$ across the three datasets are the minimum
prerequisite for a physics interpretation of the cross-dataset shape
comparison.  Non-uniform closure across datasets folds directly into
the systematic uncertainty on the ratio.

\begin{figure}[!t]
  \centering
  \includegraphics[width=\linewidth]{figs/sensitivity_ratios_placeholder}
  \caption{Cross-dataset per-event yield ratios
    $\Ntrue(\Usym)/\Ntrue(\Usym{=}18\,\text{MeV})$
    for $\Usym \in \{0, 90\}$~MeV (solid) versus the corresponding
    MC-truth ratios (dashed).  Deviation from unity in excess of the
    error band is the sensitivity signal.
    %% TODO(fig:ratio): swap in sensitivity_ratios.png.
  }
  \label{fig:ratio}
\end{figure}

\begin{table}[!t]
  \centering
  \caption{Sensitivity summary at threshold $t = 0.5$ (placeholder
    numbers from the smoke run; final values from full SLURM sweep).
    $N_{\text{ev}}$ is the per-dataset evaluation-split event count;
    closure is $\Ntrue/\Nmc$.}
  \label{tab:summary}
  \begin{tabular}{lrrrrrr}
    \toprule
    Dataset      & $\Usym$ & $N_{\text{ev}}$ & $\NperEv{\Nreco}$
                 & $\NperEv{\Ntrue}$ & $\NperEv{\Nmc}$ & closure \\
    \midrule
    \texttt{zeroSpot}    & 0   & \TODO{} & \TODO{} & \TODO{} & \TODO{} & \TODO{} \\
    \texttt{defaultSpot} & 18  & \TODO{} & \TODO{} & \TODO{} & \TODO{} & \TODO{} \\
    \texttt{bigSpot}     & 90  & \TODO{} & \TODO{} & \TODO{} & \TODO{} & \TODO{} \\
    \bottomrule
  \end{tabular}
\end{table}

\subsection{Threshold-scan robustness and operating-point choice}

Cluster candidacy is enforced by a single cut $s > t$ on the
per-cluster classifier score.  We treat $t$ as a
\emph{physics-optimized free parameter}, not a canonical value:
the value quoted in Sec.~\ref{sec:performance} at $t = 0.5$ is
reported for continuity with the analysis-and-detector-meeting
plots~\cite{Bocharnikov2026Dombay} and to allow direct comparison
with the earlier prototype baselines of Ref.~\cite{Morozov2024HGND},
but the final analysis operating point is chosen to maximise a
physics-driven figure of merit rather than to hit a fixed classifier
score.  Two conventions are quoted:
\begin{itemize}
  \item \textbf{Purity-locked benchmark.} Following
    Ref.~\cite{Morozov2024HGND}, we quote the reconstruction
    efficiency vs.\ $\Ekin$ at the threshold that yields a global
    signal purity of $\pi = 0.7$.  This is the operating point used
    for cross-checks against the rule-based cluster baseline of
    Sec.~\ref{sec:gnn:baseline}.
  \item \textbf{Sensitivity-optimal.} For the cross-dataset ratio
    plots in Fig.~\ref{fig:ratio} we scan $t \in [0.05, 0.95]$ and
    pick $t^{\star}$ that maximises the significance
    $|\Ntrue(\Usym{=}0) - \Ntrue(\Usym{=}90)|
     / \sqrt{\Ntrue(\Usym{=}18)}$
    integrated over $0.5 \le \Ekin \le 4$~GeV.  The stability of the
    resulting ratio under $t \to t \pm 0.15$ is reported as a
    threshold systematic.
\end{itemize}
Figure~\ref{fig:threshold} overlays cross-dataset ratios for
$t \in \{0.3, 0.5, 0.7\}$; a stable shape across thresholds indicates
the sensitivity signal is not an artefact of a particular operating
point.
%% TODO(sec:sensitivity): confirm the smoke-scale trend survives at
%% full scale.  If not, defer the qualitative sensitivity claim.

\begin{figure}[!t]
  \centering
  \includegraphics[width=\linewidth]{figs/sensitivity_threshold_scan_placeholder}
  \caption{Threshold-scan robustness of the cross-dataset per-event
    yield ratios.
    %% TODO(fig:threshold): swap in sensitivity_threshold_scan.png.
  }
  \label{fig:threshold}
\end{figure}

% =====================================================================
\section{Systematic uncertainties}
\label{sec:systematics}
% ---------------------------------------------------------------------
% Enumerate what has been quantified vs. what is deferred.  Anticipate
% the reviewer questions on: scaler drift, model family, seed variance,
% closure bias, MC-truth definition.
% ---------------------------------------------------------------------

We identify the following sources of systematic uncertainty on the
per-event yield ratio.  Each is discussed with the mitigation in
place at time of writing and the reference plot or table where the
impact is quantified.

\begin{enumerate}
  \item \textbf{Cross-dataset feature-scaler drift.}
    Each \SMASH{} dataset fits its own \texttt{StandardScaler} during
    preprocessing (see Sec.~\ref{sec:simulation}).  Evaluating a model
    trained on \texttt{defaultSpot} on the two other datasets
    therefore introduces a controlled input-distribution shift.
    Impact on the closure metric was tested by refitting a pooled
    \texttt{StandardScaler} on the concatenation of top and bottom
    hits from all three datasets, remapping every graph via the
    analytic $x_{\text{new}} = (\sigma_{\text{old}}/\sigma_{\text{new}})\,x_{\text{old}} + (\mu_{\text{old}} - \mu_{\text{new}})/\sigma_{\text{new}}$
    identity, and re-evaluating the HPC checkpoint on the same
    matched 4-shard subsets (see
    \texttt{scripts/run\_pooled\_scaler\_sensitivity.py}).  On the
    smoke-scale sample the closure shifts by only $+0.009$
    (\texttt{zeroSpot}), $0.000$ (\texttt{defaultSpot}), and $+0.008$
    (\texttt{bigSpot}) --- \texttt{defaultSpot} dominates the pooled
    fit by 285:1 in hit count so the pooled scaler is essentially
    identical to \texttt{defaultSpot}'s own.  The 37\,\% closure
    deficit on \texttt{defaultSpot} is therefore \emph{not} caused
    by scaler contamination; likely candidates are the
    $e_{\text{pred}}$-vs-$e_{\text{true}}$ binning mismatch and
    residual multi-neutron cluster ambiguity (Sec.~\ref{sec:sensitivity}).
    \TODO{re-test at full \SMASH{} statistics after the SLURM sweep;
    OOD scaler drift may re-emerge when the three datasets contribute
    comparable hit counts to the pool.}
  \item \textbf{Non-uniform closure.}
    A dataset-dependent $\Ntrue/\Nmc$ folds directly into the ratio
    plot; we quote the maximum $|\text{closure}(d) - 1|$ across the
    three datasets as an additional systematic on the ratio.
  \item \textbf{Model-family disagreement.}
    We reweight the analysis using \texttt{hetero\_hgt} and
    \texttt{spectral\_dynedge} checkpoints and quote the spread on
    the per-event ratio as a model systematic.
    %% TODO(sec:sys): produce a comparison table across model families.
  \item \textbf{Seed variance.}
    Three independent training seeds per model provide the intra-model
    training noise; quoted as an additional error band.
  \item \textbf{Threshold choice.}
    Included via the scan in Fig.~\ref{fig:threshold}.
  \item \textbf{Binning + MC-truth definition.}
    We report per-cluster ($\texttt{cl\_label}=1$) and per-neutron
    (unique \texttt{(Row, Id)} tracks with $\texttt{n0\_label}=1$)
    conventions consistently; downstream ratios are invariant under
    the choice provided the same convention is used in numerator and
    denominator.
  \item \textbf{Detector-response systematics.}
    Time-resolution smearing, per-cell energy calibration, and dead
    channels follow Ref.~\cite{Morozov2024HGND}; we do not repeat
    those studies here.
\end{enumerate}

% =====================================================================
\section{Summary and outlook}
\label{sec:summary}
% ---------------------------------------------------------------------

We have presented a graph-neural-network based neutron reconstruction
pipeline for the \HGND{} at \BMN{} and used it to establish the first
end-to-end sensitivity benchmark of the detector concept to the
nuclear symmetry potential in three \SMASH{} Xe+Cs simulations.
%% TODO(sec:summary): write one paragraph of headline numbers once the
%% full SLURM sweep is complete.

The next steps toward first-data readiness are:
(i) closure-driven refinement of the efficiency correction so that
$\Ntrue/\Nmc \to 1$ uniformly across datasets;
(ii) a differentiable clustering step to replace the fixed
$r_{\text{loc}}, \Delta t$ radius-graph;
(iii) a domain-adaptation study to bridge the MC$\to$real-data
gap using the first \BMN{} Xe+Cs sample;
(iv) inclusion of the proton observables from complementary \BMN{}
subdetectors to convert the neutron-only sensitivity into an
$n/p$-ratio measurement directly comparable to
FOPI/LAND~\cite{FOPI_LAND1, FOPI_LAND2} at BM@N energies.

% =====================================================================
\begin{acknowledgments}
This work was carried out at the Institute for Nuclear Research,
Russian Academy of Sciences, and supported by the Russian Scientific
Foundation grant No.~22-12-00132.  This research received funding from
the Basic Research Program at the National Research University Higher
School of Economics.  Computational resources were provided by the
HPC facilities at HSE University.
\end{acknowledgments}

% =====================================================================
\appendix
\section{Reproducibility}
\label{sec:reproducibility}
% ---------------------------------------------------------------------
% Mention the analysis repo, SLURM/local scripts, seed, dataset
% checksums.  Ties directly to notebooks/scripts/slurm layout.
% ---------------------------------------------------------------------

All analysis code, training scripts, and figure-producing notebooks
are released with this paper at
\CITENEEDED{repo url + tag}.
The pipeline supports two execution modes with identical code paths:
\begin{itemize}
  \item \emph{Local (macOS/MPS, laptop).}
    Smoke-scale reproduction runs on a subset of the \SMASH{} data
    via \texttt{scripts/run\_sensitivity\_smoke.sh} in
    $\sim 5$~min wall time; peak resident memory
    $< 3$~GB.  A memory-safe device planner
    (\texttt{HGNDRecoGNN.device}) auto-selects CUDA $\to$ MPS $\to$
    CPU and applies per-submodule CPU pinning for MPS-hostile
    ops (\texttt{DynamicEdgeConv} in particular).
  \item \emph{Cluster (SLURM/CUDA).}
    Full-scale runs are driven by
    \texttt{slurm/sensitivity\_full.sbatch} and support array-job
    sweeps across models via \texttt{--array=0-2}.  Per-run wall
    time is \NUM{6}{\text{h}}{measure} on a single A100.
\end{itemize}
Random seed 42 is used throughout; deterministic
\texttt{torch\_geometric} loaders are enabled.  Dataset shard
manifests (\texttt{meta.json}) record schema version, event count,
and per-shard checksum.

%% TODO(sec:reproducibility): add the git tag, DOI, and dataset
%% checksums once the paper is submitted.

% =====================================================================
% Bibliography
% =====================================================================
\bibliographystyle{apsrev4-2}
% Inline biblist for the sketch; migrate to main.bib before submission.
\begin{thebibliography}{99}

\bibitem{Morozov2024HGND}
  S.~Morozov \emph{et~al.},
  ``The Highly-Granular Time-of-Flight Neutron Detector for the BM@N
   experiment,''
  arXiv:2412.00455 (2024).

\bibitem{Guber2023HGND}
  F.~Guber \emph{et~al.},
  ``Development of high granular neutron time-of-flight detector for
   the BM@N experiment,''
  arXiv:2309.09610 (2023).

\bibitem{Finogeev2024TDC}
  D.~Finogeev \emph{et~al.},
  ``Development of a 100~ps TDC based on a Kintex 7 FPGA for the high
   granular neutron time-of-flight detector for the BM@N experiment,''
  Nucl.~Instrum.~Meth.~A \textbf{1059}, 168952 (2024).

\bibitem{Guber2023SiPM}
  F.~Guber \emph{et~al.},
  ``Measurement of time resolution of scintillation detectors with
   EQR-15 SiPMs for the ToF neutron detector of the BM@N experiment,''
  arXiv:2309.03614 (2023).

\bibitem{Afanasiev2024BMN}
  S.~Afanasiev \emph{et~al.},
  ``The BM@N spectrometer at the NICA accelerator complex,''
  Nucl.~Instrum.~Meth.~A \textbf{1065}, 169532 (2024);
  arXiv:2312.17573.

\bibitem{BMN_EPJ}
  M.~Kapishin (BM@N Collaboration),
  ``Studies of baryonic matter at the BM@N experiment (JINR),''
  Nucl.~Phys.~A \textbf{982}, 967 (2019).

\bibitem{Cleymans}
  J.~Randrup and J.~Cleymans,
  ``Maximum freeze-out baryon density in nuclear collisions,''
  Phys.~Rev.~C \textbf{74}, 047901 (2006).

\bibitem{Friman}
  B.~Friman, W.~N\"orenberg, and V.~D.~Toneev,
  ``The quark condensate in relativistic nucleus-nucleus collisions,''
  Eur.~Phys.~J.~A \textbf{3}, 165 (1998).

\bibitem{Senger2022BMN}
  P.~Senger (BM@N Collaboration),
  ``The heavy-ion program at the upgraded BM@N experiment at NICA,''
  PoS \textbf{CPOD2021}, 033 (2022).

\bibitem{Sorensen2024EoS}
  A.~Sorensen \emph{et~al.},
  ``Dense nuclear matter equation of state from heavy-ion collisions,''
  Prog.~Part.~Nucl.~Phys.~\textbf{134}, 104080 (2024).

\bibitem{LeFevre2015}
  A.~Le~F\`evre \emph{et~al.},
  ``Constraining the nuclear matter equation of state around twice
   saturation density,''
  Nucl.~Phys.~A \textbf{945}, 112 (2016).

\bibitem{E895_ellipticFlow}
  C.~Pinkenburg \emph{et~al.} (E895 Collaboration),
  ``Elliptic flow: transition from out-of-plane to in-plane emission
   in Au+Au collisions,''
  Phys.~Rev.~Lett.~\textbf{83}, 1295 (1999).

\bibitem{E895_sideFlow}
  H.~Liu \emph{et~al.} (E895 Collaboration),
  ``Sideward flow in Au+Au collisions between 2 and 8~AGeV,''
  Phys.~Rev.~Lett.~\textbf{84}, 5488 (2000).

\bibitem{E895_diffElliptic}
  P.~Chung \emph{et~al.} (E895 Collaboration),
  ``Differential elliptic flow in 2--6 AGeV Au+Au collisions: a new
   constraint for the nuclear equation of state,''
  Phys.~Rev.~C \textbf{66}, 021901 (2002).

\bibitem{Burgio2020NuclearMatter}
  G.~F.~Burgio and I.~Vidana,
  ``The Equation of State of Nuclear Matter: from Finite Nuclei to
   Neutron Stars,''
  Universe \textbf{6} (8), 119 (2020).

\bibitem{FOPI_LAND1}
  Y.~Leifels \emph{et~al.},
  ``Exclusive studies of neutron and charged particle emission in
   Au+Au at 400~MeV/nucleon,''
  Phys.~Rev.~Lett.~\textbf{71}, 963 (1993).

\bibitem{FOPI_LAND2}
  D.~Lambrecht \emph{et~al.},
  ``Energy dependence of collective flow of neutrons and protons in
   Au+Au collisions,''
  Z.~Phys.~A \textbf{350}, 115 (1994).

\bibitem{LongWei2022}
  X.-X.~Long and G.-F.~Wei,
  ``Effects of incompressibility $K_0$ in heavy-ion collisions at
   intermediate energies,''
  Phys.~Rev.~C \textbf{109}, 054619 (2024).

\bibitem{SMASH}
  J.~Weil \emph{et~al.} (\SMASH{} Collaboration),
  ``Particle production and equilibrium properties within a new
   hadron-transport approach for heavy-ion collisions,''
  Phys.~Rev.~C \textbf{94}, 054905 (2016);
  \CITENEEDED{SMASH DOI update if newer reference exists}.

\bibitem{Geant4}
  S.~Agostinelli \emph{et~al.} (\GEANT{} Collaboration),
  ``Geant4 --- a simulation toolkit,''
  Nucl.~Instrum.~Meth.~A \textbf{506}, 250 (2003).

\bibitem{bmnroot}
  BM@N Collaboration,
  ``\texttt{bmnroot}: simulation and analysis framework for the BM@N
   experiment,''
  \url{https://git.jinr.ru/nica/bmnroot}.

\bibitem{Baznat2020DCMSMM}
  M.~Baznat, A.~Botvina, G.~Musulmanbekov, V.~Toneev, and V.~Zhezher,
  ``Monte-Carlo Generator of Heavy Ion Collisions DCM-SMM,''
  Phys.~Part.~Nucl.~Lett.~\textbf{17}, 303 (2020).

\bibitem{Wilson1927}
  E.~B.~Wilson,
  ``Probable inference, the law of succession, and statistical inference,''
  J.~Amer.~Statist.~Assoc.~\textbf{22}, 209 (1927).

\bibitem{FeyLenssen2019PyG}
  M.~Fey and J.~E.~Lenssen,
  ``Fast Graph Representation Learning with PyTorch Geometric,''
  ICLR Workshop on Representation Learning on Graphs and Manifolds (2019).

\bibitem{KipfWelling2016GCN}
  T.~N.~Kipf and M.~Welling,
  ``Semi-supervised classification with graph convolutional networks,''
  arXiv:1609.02907 (2016).

\bibitem{Bocharnikov2026Dombay}
  V.~Bocharnikov,
  ``Graph Neural Network-based neutron reconstruction in the HGND at
   the BM@N experiment,''
  presentation at the BM@N Analysis and Detector Meeting,
  26 February 2026 (unpublished).

\end{thebibliography}

\end{document}
